\documentstyle[12pt,fullpage,trademark,epsf,apalike]{article}

%\renewcommand{\baselinestretch}{1.5}

% needed for the picTeX figures
\input prepic
\input pictex
\input postpic

%\flushbottom

\begin{document}

\renewcommand{\topfraction}{1.0}
\renewcommand{\dbltopfraction}{1.0}
\renewcommand{\bottomfraction}{1.0}
\renewcommand{\textfraction}{0.0}
\textfloatsep 0.1in
\dbltextfloatsep 0.1in

\bibliographystyle{apalike}

\nocite{As85,AB96,AB97,ADH93,AD98,BDLS96,%
BA86b,BCW88,BH92b,BC96,BCS96,Bao97,BA97,Chau82,Cha97,Cha98,%
CCMS94,CMSC96,CSMC97,Co91,Cr84,Cr93,CLNL87,COW92,%
DRMVKSD94,Dyk96,EE98,ESRI96,GHW92,GMSBK97,Gue98,%
HaBr91b,HKT95,HBCUW91,HMW96,KM94,LSDB92,%
MacD92,Math96,MCCKLS95,MCCKLSv95,MCCS96,MCSMv96,MC98,MW87,Mon88,Mon89,%
Nye92,OP96,OPR96,Pajarola:98th,Pan96,RMJF92,Sc94,St90,SCB98,%
SMCC94,SMC95,SMCv95,SMMC96Int,SMC96ACM,%
SMC97IASC,SKSKCS97Int,SKSCL97ASA,SCKL98Army,SPW98ASA,UWH90,WLB90,ZG97}

\hyphenation{Arc-View}

\title{\Large\bf Linking \ArcViewTHREEZEROTM\/ and XGobi:\\
Insight Behind the Front End}
\author{\normalsize\bf J\"urgen Symanzik$^{\normalsize 1}$, 
Dianne Cook$^{\normalsize 2}$,
Nicholas Lewin--Koh$^{\normalsize 2}$, \\
\normalsize\bf James J. Majure$^{\normalsize 2}$,
Inna Megretskaia$^{\normalsize 2}$\\
{\small {\normalsize $^1$} Utah State University, Department of
Mathematics and Statistics, Logan, UT 84322--3900, USA}\\
{\small {\normalsize $^2$} Iowa State University, 
Department of Statistics, Ames, IA 50011, USA}\\
{\small \tt symanzik@sunfs.math.usu.edu}}
\date{}
\maketitle

\vspace{-0.2in} \abstract This paper presents aspects of the
implementation of a bidirectional link between the Geographic
Information System (GIS) \ArcViewTHREEZEROTMfn\/ and the interactive
dynamic statistical graphics program XGobi. We describe the main
functionality of the link, the underlying Remote Procedure
Call (RPC) mechanism, and internal data structures, and discuss
topics such as security, concurrency, and linked brushing. We think
that these topics are of particular interest to software authors
intending to link similar software packages, and software users
learning about strengths (and weaknesses) of the implementation of our
link.

{\bf Keywords:}
Cartographic data;
Dynamic statistical graphics; 
Geographic information system; 
Interprocess communication;
Multivariate spatial data;
Spatial cumulative distribution function;
Spatial statistics;
Spatially lagged scatterplot;
Variogram--cloud plot.

\normalsize


\section{Introduction}

Geographic Information Systems (GISs) are used extensively for
examining spatial data but, traditionally, they have few facilities
for exploratory spatial data analysis. Likewise, dynamic statistical
graphics software packages are used primarily for exploratory data
analysis and have little of the sophisticated cartographic tools
necessary for displaying geographic information. The lack of
exploratory tools in a GIS is especially critical when there are
multiple attributes measured at each spatial location, and the lack of
cartographic tools in dynamic graphics software is critical in the
presence of extensive data bases of geographic information.

To gain access to sophisticated dynamic graphics tools from an
extensive GIS, we developed a bidirectional link between a GIS,
ArcView, and a dynamic graphics program, XGobi (Swayne, Cook \&
Buja, 1998). ArcView is a widely--used and established GIS.  XGobi is
a dynamic graphics program which is especially powerful for exploring
multivariate data. A seamless link is a natural way to obtain the
desired capabilities, for several reasons: (1) ArcView has an
established clientele but the graphical design makes it virtually
impossible to build in the capabilities of XGobi, (2) XGobi's visual
paradigms make it almost impossible to build in the sophisticated
mapping tools available in ArcView, and (3) we use both packages in
our own data analyses and the link is a way to streamline the sharing
of information between the two programs. Both packages are unique and
complementary. Their paths meet in the current climate of increasing
amounts of multivariate geographically referenced information.
ArcView has some rudimentary tools for establishing InterProcess
Communication (IPC) through its Remote Procedure Call (RPC) mechanism, the
basis for our link.

In general, seamless software links through IPC allow the rapid
prototyping of cross--disciplinary methodology into complex
software systems without having to re--implement existing
individual software parts. This paper describes
how RPCs were built into XGobi and the modifications made to ArcView
which allowed the two packages to share data. It will be informative for
software authors who wish to link other software packages and also to
users wishing to learn about the strengths and weaknesses of our
implementation.

The next section discusses the literature on geographic visualization
and dynamic statistical graphics.  In Section 3 we give a short
summary of IPC standards and, in particular, the RPC
mechanism. Section 4 provides an introduction to the capabilities of
the link between ArcView 3.0 and XGobi. Section 5 describes our
implementation in detail. We finish with a discussion in Section 6.

Our ArcView/XGobi link was developed on \DECTMfn\/ alphastations, and
tested on \SunTMfn/\SparcTMfn\/, \SGITMfn\/, and Data General
workstations.  A current version of the software can be downloaded from
the web at the URL 
\begin{verbatim}
    http://www.public.iastate.edu/~arcview-xgobi/homepage.html.
\end{verbatim}
A description of the software from the user's point of view is given
in Symanzik, Majure \& Cook (1997c). Various pieces of this work have been
documented in conference proceedings. This paper represents a
comprehensive description of what is required to get an RPC link
running.


\section{Cartographic and Dynamic Statistical Graphics}

Linking statistical plots with geography for
analyzing spatially referenced data has been discussed widely in
recent years. Monmonier (1988) describes a conceptual framework for
geographical representations in statistical graphics and introduces
the term {\it geographic brushing} in reference to interacting with the
map view. But geographic brushing does not only mean pure interaction
with the map. In addition, this term has a much broader meaning,
e.~g., finding neighboring points and spatial structure in a 
geographic setting.

A good overview of dynamic graphics for multivariate data is given by
Buja, Cook \& Swayne (1996).  In addition, many software solutions have
been developed for exploring multivariate spatially referenced data.
We describe a sampling of this work here.

In McDonald \& Willis (1987), a grand tour (Asimov, 1985; Buja \&
Asimov, 1986) is linked to an image to assess the clustering of
landscape types in the band space of a LandSat image taken over
Manaus, Brazil. In Carr, Littlefield,
Nicholson \& Littlefield (1987) and Monmonier (1989), a scatterplot
matrix is linked to a map view.  In REGARD (Unwin, Wills \& Haslett,
1990), map views are linked with histograms and scatterplots and,
moreover, diagnostic plots for assessing spatial dependence are also
available. Another exploratory system that links histograms and
scatterplots with latitude and longitude (and depth) coordinates is
discussed in MacDougall (1992).  In Carr, Olsen \& White (1992),
(bivariate) ray--glyph maps have been linked with scatterplots.
Klein \& Moreira (1994) report on an
interface between the image program MTID and XGobi, used for the
exploratory analysis of agricultural images.  
DiBiase, Reeves, MacEachren, von Wyss, Krygier, Sloan \& Detweiler
(1994) provide an overview on existing multivariate (statistical)
displays for geographic data.  Some recent
developments are the cartographic data visualizer, {\tt cdv} (Dykes,
1996), where a variety of plots are linked with geography, the
Space--Time--Attribute Creature/Movie, STAC/M (Openshaw \& Perr{\'e}e,
1996), that searches for patterns in GIS data bases under the control
of a Genetic Algorithm, and an exploratory spatial analysis system in
XLisp--Stat (Brunsdon \& Charlton, 1996).  The reader of these
references should note that in most geographically influenced
publications authors distinguish between {\it geographic (or
cartographic) space} and {\it attribute (or data) space} but rarely
use the statistical expression {\it variable} to relate to the latter
one.

In addition to the ArcView/XGobi link, there are several other examples
where GISs and (graphical) statistical packages have been linked.
Williams, Limp \& Briuer (1990) demonstrate how S and the GRASS GIS
can be jointly used for archaeological site classification and analysis.
Scott (1994) links STATA with ArcView and the spatial data analysis
software SpaceStat has been linked with \ARCINFOTMfn\/ (Anselin,
Dodson \& Hudak, 1993) and with ArcView 
(Anselin \& Bao, 1996; Anselin \& Bao, 1997). 
In Haining, Ma \& Wise (1996), the designing of a software system for
interactive exploration of spatial data by linking to ARC/INFO has
been discussed, and in Zhang \& Griffith (1997), a spatial statistical
analysis module implemented in ArcView using Avenue has been discussed. 
Mathsoft (1996) describes the S+GISLink, a
bidirectional link between ARC/INFO and \SPLUSRfn, and Bao (1997)
describes the S+Grassland link between S--PLUS and the Grassland
GIS. Finally, a comparison of the operational issues of the
SpaceStat/ArcView link and the S+Grassland link has been given in Bao
\& Anselin (1997).


\section{InterProcess Communication Standards and Remote Procedure Calls}

The main goal of our research was to establish a bidirectional link
between ArcView (first versions 2.0 and 2.1, now version 3.0) and
XGobi. This link utilizes Remote Procedure Calls (RPCs), an
InterProcess Communication (IPC) feature available in ArcView.  The
use of RPCs is a programming technique where a process on the local
system ({\em client}) invokes a procedure on a remote system ({\em
server}).  In this context, the term {\em request} is used to refer to
the client's desire to execute a remote procedure and the
term {\em response} is used for the result produced by the remote
procedure (Stevens, 1990).  
RPCs are differently implemented on different hardware platforms
(see Stevens, 1990, Chapter 18, for more details).
It should be possible to use the ArcView/XGobi link on any
hardware platform that supports Sun Microsystems' RPC mechanism
in addition to DEC alpha\-stations and
Sun/Sparc, SGI, and Data General workstations
that all conform to this RPC standard.
Detailed information on this RPC mechanism can be found
in Corbin (1991). Stevens (1990) describes other IPC standards such as
pipes, FIFOs (named pipes), message queues, semaphores, shared memory,
Berkeley sockets, and the System V Transport Layer Interface (TLI) IPC
on \UnixRfn\/ systems.

RPCs are a mature and robust IPC technology easily 
applied between ArcView and XGobi. Newer distributed software
tools like the Parallel Virtual Machine (PVM, e.~g., Bode, Dongarra,
Ludwig \& Sunderam, 1996; Alexandrov \& Dongarra, 1998),
the Distributed Component Object Model (DCOM, e.~g., Eddon \& Eddon, 1998),
and the Common Object Request Broker Architecture (CORBA, e.~g.,
Otte, Patrick \& Roy, 1996) have not been explored for use in the
ArcView/XGobi link.
The reason is easy. In 1994, when we started work on this link,
ArcView 2.0 supported only RPCs as IPC mechanism. Even now,
at publication date,
it would be extremely difficult to implement any additional IPC
technology in ArcView without having access to the ArcView source code.

For linking geographic packages with statistical packages there are
two existing approaches: {\it close coupling} and {\it loose coupling}
(Goodchild, Haining \& Wise, 1992). Close coupling relates to the fact
that one package is calling other packages directly or that
this package is calling user--written
routines using one of the previously listed IPC facilities. Loose coupling is
simpler: multiple processes simultaneously access (read and write) the
same file(s) or they simply pass data files (ASCII or binary) among each
other. The link between ArcView and Space\-Stat (Anselin \& Bao, 1996;
Anselin \& Bao, 1997)
and the software system designed by Haining et al.\ (1996) are
examples for loose coupling.  The S+Grassland link (Bao, 1997) and the
ArcView/XGobi link are examples of close coupling.

Overall, the approach of linking existing statistical and geographic
packages through close coupling has several advantages when compared
to linking through loose coupling.  Close coupling in general provides
access to the full capabilities of both packages without
re--developing sophisticated existing software. It represents the
current trend towards {\it open systems} (G{\"u}nther, 1998, p.\ 143).
Moreover, linked brushing, and truly dynamic graphics capabilities, is
practically feasible only through close coupling, and it considerably
enhances the power of the software environment beyond the sum of the
capabilities of the individual packages.

IPC facilities are likely to be incorporated in more
statistical software packages over time to make them accessible in
open systems. As an example, the latest revision of S, ``Version 4'',
provides a so--called {\it connection class object} which has specific
classes such as files, pipes, and FIFOs.  The connection class can be
extended towards other IPC connections as well (Chambers, 1997;
Chambers, 1998).
Another similar development that will likely gain more importance in
the near future is the linking of statistical and non--statistical
software that is accessible through the World Wide Web.  The Method
ManageMent system MMM (G\"unther, M\"uller, Schmidt, Bhargava \&
Krishnan, 1997) is such a web--based system that currently links three
software packages, i.~e., Mathematica, Matlab, and XploRe.
``The Omega Project for Statistical Computing''
(\verb|http://www.omegahat.org/|)
has the goal to
provide open--source software for statistical applications and support for
distributed statistical computing through the use of the 
CORBA (e.~g., Otte, Patrick \& Roy, 1996) standard. 



\section{The Capabilities of the ArcView/XGobi Link}

After experimenting with linking a GIS and a statistical
package by implementing a unidirectional link between ARC/INFO
and XGobi in 1993 and early 1994 (Cook, Cressie, Majure \& Symanzik, 1994;
Symanzik, Majure, Cook \& Cressie, 1994), we restarted the
work using ArcView 2.0. 
ArcView 2.0 supported calls to external procedures
via RPCs and enabled external programs to invoke
ArcView 2.0 functions using RPCs in contrast to ARC/INFO. 
These features enabled us
to implement the kernel for a bidirectional link between
ArcView 2.0 and XGobi that is still central to the current 
version of the link.

Initially, we concentrated on one--to--one connections between
ArcView and XGobi such as linking geographic location to (multivariate)
attribute values and to empirical
spatial cumulative distribution function (SCDF) values.
Subsequently, we explored one--to--two and two--to--one linking
to connect variogram--cloud plots (Chauvet, 1982; Cressie, 1984;
Haslett, Bradley, Craig, Unwin \& Wills, 1991; Haslett \& Bradley,
1991; Bradley \& Haslett, 1992), spatially lagged scatterplots
(Cressie, 1984; Rossi, Mulla, Journel \& Franz, 1992), and
multivariate variogram--cloud plots (Majure \& Cressie, 1998) 
to geography. 

We call these five diagnostic plots {\it links} as well when used
in the ArcView/XGobi context.
When speaking of the {\it basic (or multivariate) link},
we refer to the feature of the ArcView/XGobi link that links 
geographic location to (multivariate) attribute values.
Alternatively,
we refer to the other features as
{\it SCDF link}, {\it variogram--cloud link}, 
{\it spatially lagged scatterplot link},
and {\it multivariate variogram--cloud link}. This dual use
of the term {\it link} is intentional. Although these five
features are a built--in part of XGobi, they cannot be used
as a stand--alone part within XGobi, i.~e., they cannot be invoked
through a menu option in XGobi. 
When using the term {\it link} in this paper, we refer to
all five features of the ArcView/XGobi link and their
underlying implementation using RPCs.

The linking between ArcView and XGobi allows us to simultaneously display 
spatial locations and concomitant
geographic variables within the GIS while visualizing and exploring
the corresponding data space within XGobi. The
usefulness of the link has been highlighted for several different
applications such as satellite imagery, forest health monitoring, and
precipitation data (Cook, Majure, Symanzik \& Cressie, 1996; Majure,
Cook, Cressie, Kaiser, Lahiri \& Symanzik, 1996a; Majure, Cressie,
Cook \& Symanzik, 1996c; Symanzik, Majure \& Cook, 1996; Symanzik,
Majure \& Cook, 1997b; Symanzik, Megretskaia, Majure \& Cook, 1997d).  

Cook, Symanzik, Majure \& Cressie (1997)
presents examples (measures of livability in the United States, forest
health data, and precipitation measurements) for all five features of
the link. There also exist videos that demonstrate the use of the link
(Majure, Cook, Cressie, Kaiser, Lahiri \& Symanzik, 1995; Symanzik,
Majure \& Cook, 1995; Majure, Cook, Symanzik \& Megretskaia, 1996b).

In the remainder of this section, we explain the general
idea of linked brushing between ArcView and XGobi and
provide two examples for the
multivariate link and the variogram--cloud link.


\subsection{Linked Brushing between ArcView and XGobi}

In the multivariate and SCDF links, there is exactly one point in XGobi
corresponding to exactly one spatial location in ArcView 
and vice versa. For each point/location brushed, 
the corresponding location/point is brushed
as well.  In contrast, for the variogram--cloud, spatially lagged
scatterplot, and multivariate variogram--cloud links, every point in
XGobi corresponds to exactly two locations in ArcView. Actually, there
might be up to $n^2$ points in XGobi if there are $n$ locations in
ArcView.  We assume we brush only one point/location at a time in either
application.  Then, we make use of the following strategy for linked
brushing:

\begin{description}
\item[Brushing in XGobi:] First we mark the point in XGobi with the 
selected color, glyph, and size. Then we determine in XGobi which
are the two related spatial locations (using {\it Array of Pairs}
described in Section 5.2.2)
and pass this information to ArcView. Finally, we 
connect these locations in ArcView with a line of the
selected color (provided by XGobi). If we brush a point in XGobi
that relates to locations with identical spatial origin
($\mbox{\boldmath$s_i$} = \mbox{\boldmath$s_j$}$) in ArcView,
we do not plot anything on the ArcView side. 

\item[Brushing in ArcView:] This geographic brushing step
can be described as a three phase
process, where ArcView, XGobi, and again ArcView proceed. 
First, we brush the location in the ArcView map view (nothing visible
happens) and pass the information to XGobi, which location has been brushed
and which color, glyph, and size have been selected.
Then, in XGobi, we determine all points that are related
to this location (using {\it Brushing List 1} and {\it 2} 
described in Section 5.2.2).
There might be up to $2 n - 1$ such points in XGobi. For each
of these points, we assume it actually has been brushed in XGobi.
We follow the instructions on brushing in XGobi, using the color,
glyph, and size provided by ArcView. Finally, this will 
result in lines drawn in ArcView  between the selected location and all
locations that are not farther away than the cutoff distance $d_{cut}$.
\end{description}

It is important to note that in our application we redraw the
selected theme in ArcView as graphics 
with the list of coordinates stored in a
stack accessed by a hash function (a data dictionary in Avenue). 
After the data is initially passed from ArcView into XGobi,
all brushing in ArcView occurs by manipulating graphics on the screen 
directly without accessing ArcView's theme table any more.


\begin{figure*}[b] 
\centerline{\epsfxsize=6in\epsfbox{ex_multi.ps}} 
\caption{Example for the basic (or multivariate) link. The attribute
data is displayed in XGobi (left) with a full complement of
multivariate graphical tools available. This data is linked by
brushing to the map view provided by ArcView (right). In the XGobi plot, two
extremes on third quarter precipitation and elevation are painted as
solid rectangles and they are shown to be located very close
to one another in New Hampshire near the border with Maine
in the ArcView map view.}
\label{ex_mult.ps}
\end{figure*}


\subsection{Two Motivational Examples}

We describe two examples which cover two of the five main features
of the ArcView/XGobi link: the basic (or multivariate) link
and the variogram--cloud link. The data in each case shows
precipitation recordings in the northeastern United States, i.~e., in
Maine, New Hampshire, Vermont, and Massachusetts. Figure
\ref{ex_mult.ps} shows an example of the basic (or multivariate) link,
where all the attribute information is passed from ArcView
into XGobi. The example
shows the detection of outliers in the multivariate data space. 

\begin{figure*}[b]
\centerline{\epsfxsize=5in\epsfbox{ex_vario.ps}}
\caption{Example for the variogram--cloud link.
In the upper left XGobi, we have brushed
(using a solid rectangle) the highest values in the 
variogram--cloud plot ($g_{11} (\mbox{\boldmath$s_i$}, 
\mbox{\boldmath$s_j$})$) versus
$d (\mbox{\boldmath$s_i$}, \mbox{\boldmath$s_j$})$).
In the lower right ArcView map view, each pair
of locations, related to a point that has been brushed in XGobi,
has been connected by a line.}
\label{ex_vario.ps}
\end{figure*}

A trickier problem is the detection of spatial outliers, i.~e., sites that
have dramatically different attribute values from their close
neighbors, but are not noticeably different in the overall range of the
variables.  A common approach to detecting spatial outliers is to use
a variogram--cloud plot due to Cressie (1984). Here, $g_{11}
(\mbox{\boldmath$s_i$}, \mbox{\boldmath$s_j$}) = \mid Z_1
(\mbox{\boldmath$s_i$}) - Z_1 (\mbox{\boldmath$s_j$}) \mid^{1/2}$ is
plotted versus the Euclidean distance between the locations, $d
(\mbox{\boldmath$s_i$}, \mbox{\boldmath$s_j$}) = \mid\mid
\mbox{\boldmath$s_i$} - \mbox{\boldmath$s_j$} \mid\mid$, where $Z_1
(\mbox{\boldmath$s_i$})$ represents the amount of precipitation
measured at location $\mbox{\boldmath$s_i$}$ during the third quarter
of 1990. This approach is demonstrated in Figure
\ref{ex_vario.ps}. Each point in the variogram--cloud corresponds to
two sampling sites in the map. High values of $g_{11}$ (the
precipitation difference between the sampling sites is large) are
brushed in the upper left (XGobi) variogram--cloud plot and the
corresponding pairs of sampling sites are identified by connecting
lines in the lower right (ArcView) map view. Locations that have
several lines connecting them to other sites are spatial anomalies,
i.~e., sites that are quite different
from their neighbors. This display shows
two such sites --- one in Maine and one in New Hampshire. None of them
was detected in the multivariate link because neither is an outlier
in the global sense.

Note that a variogram--cloud plot is a way to explore
spatial dependence and is usually the first
step of estimating a variogram (e.~g., Cressie, 1993). 
There are also tools for smoothing available in XGobi so 
the user can interactively smooth the variogram--cloud,
giving the analyst
a preliminary non--parametric estimate of the variogram.


\section{Technical Implementation Details}

This section describes the use of RPCs in the ArcView/XGobi link,
including discussion on the
additions to XGobi, modifications to ArcView, 
and security and concurrency issues.

\subsection{RPCs in the ArcView/XGobi Link}

From within ArcView, external remote procedures can be called and
external programs can invoke internal ArcView functions via RPCs.  To
utilize these features, we initially wrote intermediate C code
containing the RPCs and called XGobi as a subroutine, following the
examples provided by the XGvis software system (Littman, Swayne, Dean
\& Buja, 1992) and the interface between the image program MTID and
XGobi (Klein \& Moreira, 1994).  With the port to ArcView 3.0, the RPC
mechanism became a permanent part of the XGobi source code.

ArcView has been modified for this application using its built--in
Avenue programming language. All of the default ArcView
functionality is available, with the addition of 
several operations that are necessary to handle
the link. Specifically, ArcView has been modified to do the
following: initiate an RPC server and client, initiate and pass data
to the XGobi process, brush locations in the map view 
and request XGobi to brush the
corresponding points, and process requests from XGobi to brush locations.
Pseudo--code for ArcView and XGobi that describes this functionality in
detail can be found in Symanzik et al.\ (1996, 1997b).

Both, ArcView and XGobi, have been set up as a server as well as a client. 
ArcView (as a client) can request action from XGobi (as the server).   
A procedure in the main event loop of XGobi has to be used to check
for incoming RPC requests from clients (namely ArcView).
Note that it is not possible to give XGobi the full functionality of an
RPC server, which would require that it automatically processes all
requests from its clients. Such an automation would never
terminate and thus, no other user action within XGobi could be
processed.
Conversely, XGobi (as a client) sends requests to ArcView (as the
server), asking ArcView to update its map view in accordance with
the brushing and subsetting information of points done in XGobi. 
In the next section, we will describe
how XGobi remote procedures can be called from within ArcView.  


\subsection{XGobi Additions}

\subsubsection{XGobi's RPC Structure}

It was necessary to add 34 remote procedures to XGobi, that are
executable through the XGobi RPC server, to obtain the functionality
for the five different links.  These remote procedures
can be grouped into 6 classes: general use
(function numbers 01--03), multivariate link (11--16), SCDF
link (21--29), variogram--cloud link (31--35), spatially lagged
scatterplot link (41--46), and multivariate variogram--cloud link 
(51--55).  A client needs to know the function numbers.
It does not use
the full internal names. The full list of remote procedures can be
found as an electronic appendix to this paper at 
\begin{verbatim}
   http://www.math.usu.edu/~symanzik/papers/2000_jcgs/functdesc.ps.
\end{verbatim}
This list should also allow other software developers to establish
RPC links to XGobi from their own programs. 

We will describe the working of the multivariate link. To activate
this link, we first call
\verb#RPC_Init_Data#(11), through which we pass the name of the data
set, the number of columns (ncols), and the number of rows (nrows)
from ArcView into XGobi. Many XGobi data structures are initialized
here, internal memory to store the data (which is sent in a later
call) is allocated, and the flag \verb#is_init# is set. We call
\verb#RPC_Send_Colnames#(01) to pass the strings containing the names
of the variables into XGobi. No flags are modified. If this remote
procedure is not called, default variable names will be assigned. Also
optionally, names for each observation can be passed into XGobi
through \verb#RPC_Send_Rownames#(02). The required next steps are
\verb#RPC_Send_Init_Data#(12), \verb#RPC_Send_Init_Symbols#(13),
and \verb#RPC_Make_XGobi#(14), which sends the data matrix, the
color/glyph of each point and lastly pops up the XGobi window,
respectively. Linked
brushing in the multivariate link is now supported
through \verb#RPC_Update_All_Symbols#(15) and
\verb#RPC_Update_Some_Symbols#(16).

The remote procedure calls return information to ArcView 
for error checking. This data
consists of 5 character symbols: two digits represent the function
number, a dot serves as a separator, and another two digits describe the
result (error). ArcView uses a lookup--table of errors contained
in the file {\tt AV2XGErrTable} to provide feedback when
problems occur. This file
contains three columns of information. The first column
contains the ``{\it function number.error digits}''.  The second
column contains the severity of the error. A ``0'' means success and
is used only in combination with the error digits ``00''.  An ``R''
indicates that incorrect numerical values have been received and
ArcView should repeat (this is what the ``R'' stands for) its last
request using correct data.  A ``W'' indicates a warning. An ``F''
indicates a fatal error which is impossible to recover from so
the link should be terminated --- if possible directly from within
ArcView or through user intervention.
The third column contains the name of the remote procedure and a
verbal description of the error. 
Both, name and description, are  displayed in an ArcView
message window in case of an ``R'', ``W'', or ``F'' condition.

The full {\tt AV2XGErrTable} can be found as an electronic appendix to
this paper at
\begin{verbatim}
   http://www.math.usu.edu/~symanzik/papers/2000_jcgs/AV2XGErrTable.
\end{verbatim}
To explain how it is used, we use a small subset
of this file, i.~e., some of the lines used by the multivariate link:

\footnotesize
\baselineskip10pt
\begin{verbatim}
01.00 0 RPC_Send_Colnames: OK.
01.01 F RPC_Send_Colnames: Call RPC_Init_Data or RPC_Init_CDF_Data first.
01.02 F RPC_Send_Colnames: Submit number of columns many column names 
  (separated by space).

11.00 0 RPC_Init_Data: OK.
11.01 F RPC_Init_Data: Submit a name, number of columns, and number of rows 
  (separated by space).

12.00 0 RPC_Send_Init_Data: OK.
12.01 F RPC_Send_Init_Data: Call RPC_Init_Data first.
12.02 F RPC_Send_Init_Data: Submit (number of rows x number of columns) 
  many values (separated by space).

16.00 0 RPC_Update_Some_Symbols: OK.
16.01 F RPC_Update_Some_Symbols: Call RPC_Send_Init_Symbols and 
  RPC_Make_XGobi first.
16.02 F RPC_Update_Some_Symbols: Submit a list of number:color/glyph/size 
  (separated by space).
16.03 R RPC_Update_Some_Symbols: Wrong value for number.
16.04 R RPC_Update_Some_Symbols: Wrong value for color/glyph/size.
\end{verbatim}
\normalsize

\noindent For remote procedure 01 (\verb#RPC_Send_Colnames#) the
following problems can occur: (1) In case of error digits ``00'',
everything worked OK. (2) For error digits ``01'', 
this remote procedure has been
called before either remote procedure 11 or 21 has been called. (3) In
case of error digits ``02'', too few column names representing the
variables have been sent.  Both cases are assumed to be fatal errors
where the link should be terminated. 

In the last few lines of the lookup--table reproduced above
that relate to remote procedure 16 
(\verb#RPC_Update_Some_Symbols#), we see the use of the ``R'' condition.
This tells ArcView to repeat its last request using correct data
since an incorrect point identifier (error digits ``03'') or an
incorrect value for the encoded color, glyph, and size (error digits
``04'') has been received. Even though this particular brush request
failed, the link should still be operational and recovery from this
error seems possible, most likely just by resending the data.


\subsubsection{Modifications to XGobi's Data Structures}

This section concentrates on XGobi's 
data structures specifically created to facilitate
the variogram--cloud link, the spatially lagged scatterplot link, and
the multivariate variogram--cloud link.
There were no modifications to XGobi's internal data
structures for the multivariate link and SCDF
link. RPCs also provide access to XGobi's existing data
structures: the data matrix and the color and brush vectors. 

\input data_struct

We will explain the new data structures based on
the variogram--cloud link.  This link displays the variogram--cloud plot
and cross variogram--cloud plot in XGobi.
Originally, we have an array of {\it Initial Data} passed from ArcView
to XGobi.  See Figure \ref{DataStructs} for illustration purposes. We
will take the number of spatial locations, $n$, 
to be 3, $\mbox{\boldmath$s_i$} =
(x_i, y_i),~~i=1,\ldots,n,$ are these locations in the
2--dimensional plane, and $Z_1$ and $Z_2$ (in the general case $Z_1,
\ldots, Z_k$) are the measures taken at location $\mbox{\boldmath$s_i$}$.
Note that $Z_1$ and $Z_2$ may be the same variable taken at different
times, or they may be different variables.  
The cutoff distance, $d_{cut}$ (1.1 in this example), 
is also passed from ArcView to XGobi. 
The link will provide the user with an estimate of the expected
number of points in XGobi for a given $d_{cut}$. The larger the
number of points displayed in XGobi, the slower the link is.
However, performance considerably differs among hardware
platforms. While old Sun/Sparc workstations needed more
than 20 minutes to fully initialize the link in a complex
scenario, the same scenario has been initialized in less than
a minute on a SGI/Onyx workstation.
After some experimentation with the software, users
should know how many points can be handled reasonably well
using their computer hardware.

The {\it Array of Pairs} is the main data structure used to support
this part of the link. It 
consists of records corresponding to each pair 
$(\mbox{\boldmath$s_i$},\mbox{\boldmath$s_j$})$ of locations 
satisfying the condition
$\sqrt{(x_i - x_j)^2 + (y_i - y_j)^2} \leq d_{cut}$.
Each record consists of the location indices ($i$ and $j$), 
the angle between $\mbox{\boldmath$s_j$}-\mbox{\boldmath$s_i$}$
and the horizontal axis measured in degrees ($A$), the set 
$\gamma_{lm}, l, m = 1,\ldots,k$, 
where $\gamma_{lm} (\mbox{\boldmath$s_i$}, \mbox{\boldmath$s_j$}) = $
%
$\mid (Z_l(\mbox{\boldmath$s_i$}) - Z_m(\mbox{\boldmath$s_j$})) 
  \mid^{1/2}$, and
$D (\mbox{\boldmath$s_i$}, \mbox{\boldmath$s_j$}) = 
\sqrt{\sum_{p=1}^k (Z_p(\mbox{\boldmath$s_i$}) - 
  Z_p(\mbox{\boldmath$s_j$}))^2}$. 

With these definitions in mind, a distinction among the terms
variogram--cloud, cross variogram--cloud, and multivariate
variogram--cloud can be expressed as follows: {\it Variogram--clouds}
use $\gamma_{ii}$, {\it cross variogram--clouds} use
$\gamma_{ij}$, $i \neq j$, and {\it multivariate variogram--clouds}
use $\gamma_{ij}$ back to back with $\gamma_{ji}$, $i \neq
j$. The latter name intrinsically contains the word ``cross'' since
$\gamma_{ii}$ back to back with $\gamma_{ii}$ would be nothing else
but a mirror image of the positive $x$--axis mapped onto the negative
$x$--axis.

In many situations, we will consider a plot of $\gamma_{lm}
(\mbox{\boldmath$s_i$}, \mbox{\boldmath$s_j$})$ versus $d
(\mbox{\boldmath$s_i$}, \mbox{\boldmath$s_j$}) =$   $\mid\mid
\mbox{\boldmath$s_i$} - \mbox{\boldmath$s_j$} \mid\mid = \sqrt{(x_i -
x_j)^2 + (y_i - y_j)^2}$ (given in Figure \ref{ex_vario.ps}) to assess
the spatial dependence before fitting a (cross--) variogram model. The
corresponding variable names that appear in XGobi are called $g_{lm}$
and $d$.

The data record in XGobi's {\it Array of Pairs}
also contains $A$ (angular distance),
$\sin(A)$ and $\cos(A)$. In the geographic context, there often exists
a directional effect in the spatial dependence. A location
\mbox{\boldmath$s_1$} east of location \mbox{\boldmath$s_0$} may have
a considerable correlation with the values observed at location
\mbox{\boldmath$s_0$}, but a location \mbox{\boldmath$s_2$} at the
same distance north of location \mbox{\boldmath$s_0$} may not be
correlated at all. Note that we handle indeterminate cases
specially: When locations are identical ($i=j$), we assign $A=-90$. For
observations from the same location that are accessed through
different indices ($\mbox{\boldmath$s_i$} =
\mbox{\boldmath$s_j$}, i \neq j$), we assign $A=-45$.  So--called
``directional variograms'' are a feature of many software tools for
spatial data analysis, such as Variowin (Pannatier, 1996).  In the
ArcView/XGobi link, we can interactively overlay directional
variogram--cloud plots onto the overall variogram--cloud plot by
brushing either in a dotplot of the angle $A$ or in a scatterplot of
$\sin(A)$ versus $\cos(A)$. For example, if we want to highlight the
east--west direction, we might brush angles within the intervals
(350$^\circ$, 10$^\circ$) and (170$^\circ$, 190$^\circ$).  For the
north--south direction, we might brush angles within the intervals
(80$^\circ$, 100$^\circ$) and (260$^\circ$, 280$^\circ$).  Because
angle is a modular measurement, it is particularly convenient to
display the scatterplot of sin(A) vs cos(A) and brush ``around the
circle''. It is also more intuitive because this type of display can
be read like a compass.  Examples that feature the brushing of two
angle classes can be found in Symanzik et al.\ (1997c, 1997d) and Cook
et al.\ (1997).

A {\it Binary Map} is
created upon initialization of the link.  It is a 2--dimensional array
of size $n \cdot \mbox{ceil}(n/32)$, which is large enough to provide
information on all $n^2$ pairs of locations but small
in memory size to faciliate speed and save storage.
 A ``1'' at the $j$th bit in the
$i$th row indicates that the distance between $\mbox{\boldmath$s_i$}$
and $\mbox{\boldmath$s_j$}$ is smaller than $d_{cut}$.
For example, in the first row, third column, the ``1'' indicates
that the distance between 
$\mbox{\boldmath$s_1$}$ and $\mbox{\boldmath$s_3$}$ is smaller than
$d_{cut} = 1.1$. Using the {\it Initial Data} array, we could see that the
distance between $\mbox{\boldmath$s_1$}$ and $\mbox{\boldmath$s_3$}$ is 1.
A ``0'' indicates that the distance between 
two locations is bigger than $d_{cut}$.
For example, in the first row, second column, the ``0''
indicates that the distance between $\mbox{\boldmath$s_1$}$ and
$\mbox{\boldmath$s_2$}$ is bigger than $d_{cut} = 1.1$. Actually, it is 
$\sqrt{2}$ based on the values in the {\it Initial Data} array.
Information on about 32 pairs of locations is
packed into one unsigned integer.  Simultaneously with the creation of
{\it Binary Map}, we count the number of locations that lie within the
cutoff distance for every initial location.  After we finish looking at
{\it Initial Data}, we know exactly how many pairs that contain location
$\mbox{\boldmath$s_i$}$ as a component
are going to be in {\it Array of Pairs} for
every $i=1,\ldots,n$. This allows us to allocate memory for {\it Array
of Pairs} and {\it Brushing List 1} and {\it 2}. 
{\it Brushing List 1} and {\it 2} contain information on linking from
{\it Array of Pairs} to {\it Initial Data}. Each row 
of {\it Brushing List 1} and {\it 2} relates
to the corresponding row (and, thus, to a spatial location) in 
{\it Initial Data}, i.~e., row 1 relates to $\mbox{\boldmath$s_1$}$,
row 2 relates to $\mbox{\boldmath$s_2$}$, and
row 3 relates to $\mbox{\boldmath$s_3$}$.
{\it Brushing List 1} corresponds to all pairs of locations
for which $\mbox{\boldmath$s_i$}$ is the
first component in the pair. For example, its first row says that
location $\mbox{\boldmath$s_1$}$ is the first component in the pairs
(i.~e., rows) {\it 1} and {\it 2} in {\it Array of Pairs}.
{\it Brushing List 2} corresponds to all pairs of
locations for which $\mbox{\boldmath$s_i$}$ is the second component in the
pair. For example, its first row says that
location $\mbox{\boldmath$s_1$}$ is the second component in the pairs
(i.~e., rows) {\it 1} and {\it 5} in {\it Array of Pairs}.

Using {\it Binary Map} we pick qualifying pairs of locations, add information
to {\it Brushing List 1} and {\it 2}, and fill the record for this
pair in {\it Array of Pairs}. 
The initialization of {\it Brushing List 1} and {\it 2} and
{\it Array of Pairs} works as follows: 
For {\it Brushing List 1} and
{\it 2}, we know at each step exactly how 
many ``cells'' have been filled so far
(there exists a special array of counters). We place the
index number of the record used in {\it Array of Pairs}
in the first free positions of {\it Brushing List 1} 
and {\it 2} with respect to the two components of this pair
and increment the corresponding counters by one.  After
we have gone through {\it Binary Map}, we have {\it Array of Pairs}
and {\it Brushing List 1} and {\it 2} completely filled. 

Using these data structures, we can
brush points/locations in two directions. 
If we brush point $i$ in XGobi, we have to look at row $i$ of
{\it Array of Pairs} and know immediately which two spatial locations 
this entry consists of. Based on this information, we can 
draw a line that connects these two locations in the ArcView map view.
Alternatively, brushing location $\mbox{\boldmath$s_i$}$
in ArcView corresponds to selecting row $i$ of
{\it Initial Data}. Also,
the $i$th row of {\it Brushing List 1 (2)} contains all index numbers of
the pairs that contain location $\mbox{\boldmath$s_i$}$
as first (second) component. Since these index numbers relate
to rows in {\it Array of Pairs}, we know which points have to be 
brushed in XGobi. In continuation of our example, if we
brush location $\mbox{\boldmath$s_1$}$ in ArcView, this
translates to row 1 of
{\it Initial Data}, and, finally, points related to pairs {\it 1}, {\it 2},
and {\it 5} in {\it Array of Pairs} will be brushed in XGobi. 
See Section 4.1 for an
explanation of the general idea of linked brushing between
ArcView and XGobi.

The efficiency of our data structures is worth discussing. 
As an example, we assume there are $n = 128$ locations in ArcView initially. 
This would require an array of size $128^2 = 16384$ in XGobi if
all possible pairs of locations would be used.  
We also assume that a 
cutoff distance $d_{cut}$ has been selected such that only 1024, i.~e.,
1/16th of the possible number of points, are displayed in XGobi.
Using the {\it Binary Map} approach, we reduce the total
data array size to 1024 entries and need only 
$128^2 / 32 = 512$ unsigned integers as 
additional memory for {\it Binary Map} and 2 $\times$ 1024
entries for {\it Brushing List 1} and {\it 2}. 
This can be neglected compared to a potential
allocation of 16384 records of {\it Array of Pairs}. 
Each row of {\it Brushing List 1} and {\it 2} is reasonably
small ($1024 / 128 = 8$ in average) to allow efficient handling and
brushing of points.


\subsection{ArcView 3.0 Modifications}

The implementation of the link in ArcView was accomplished
by a combination of Avenue scripts, additions to the
graphical user interface (GUI), and, in the latest version, the
incorporation of all scripts into an extension that can be loaded by
the user at runtime. In the context of the 
ArcView/XGobi link,
ArcView has to do two things: (1) respond 
to user actions and request appropriate changes in XGobi,
and (2) respond to RPC requests from XGobi and make the requested
changes in the map view.

Figure \ref{GUI3.0.ps} shows the {\it View} Document GUI used for the
basic (or multivariate) link, the variogram--cloud link, the spatially
lagged scatterplot link, and the multivariate variogram--cloud
link. A slightly different GUI is used for the
SCDF link.

\begin{figure*}[ht]
%\vspace{1.8in}
%\special{psfile=newgui2.ps vscale=75 hscale=75 voffset=-250 hoffset=60}
\centerline{\epsfxsize=6in\epsfbox{newgui2.ps}}
% -350 far too high, -400 still too high
% have to modify hoffset: -25 shifts image to beginning of left column
%\vspace{2.0in}
\caption{The {\it View} Document GUI in ArcView 3.0,
used for all features of the link except the SCDF link. 
The following new features of the GUI have been added for the
ArcView/XGobi link but are not part of the standard
ArcView GUI. The red and green
diamond buttons allow the user to start and stop
XGobi. The paint brush button switches user control to brushing mode.
The ``XGobi'' pulldown menu allows the user to select one of the features
of the link. The
``Color'',  ``Glyph'', and ``Size'' pulldown menus allow the user to select
the plotting symbol.}
\label{GUI3.0.ps}
\end{figure*}

Features on the GUI added to facilitate the ArcView/XGobi link are:
\begin{itemize}
\item the ``XGobi'' pulldown menu:
select one of the features of the link;
\item the green and red diamond buttons: start and stop XGobi;
\item ``Color'', ``Glyph'', and ``Size'' pulldown menus:
set the current brushing characteristics, i.~e., 
color, glyph, and size of the symbol used for plotting the locations; and
\item the paint brush button: brush locations in the ArcView map view.
\end{itemize}

The following Avenue scripts, listed in alphabetical order,
are required in ArcView to support the basic 
(or multivariate) link:
\begin{description}
\item[CreateErrorDictionary:] Builds the 
table of error messages that is internally used in ArcView
based on the file {\tt AV2XGErrTable} 
(described in Section 5.2.1).
\item[RGBColorNameDictionary:] Creates a dictionary containing
ASCII color names and RGB color values that are related to the
colors used in XGobi.
\item[SetCurrentColor:] Selects and sets a 
new color in the ``Color'' pulldown menu.
\item[SetCurrentGlyph:] Selects and sets a 
new glyph in the ``Glyph'' pulldown menu.
\item[SetCurrentSize:] Selects and sets a 
new size in the ``Size'' pulldown menu.
\item[XG.Finish:] Sends a ``finish'' request to XGobi, terminates the
RPC Server within ArcView, and resets the buttons in the ArcView GUI.
\item[XG.ReadXGobiColors:] Attempts to read 
the {\tt \$HOME/app-defaults/XGobi} file (if available) and 
tries to get the current XGobi colors from it.
\item[XG.RPC.Brush.Some:] Called from within XGobi
when points are brushed in the XGobi window and an update of
the ArcView map view is required. The data that is passed is of the type
``list of index:size\_color\_glyph'' where ``index'' is the internal index
number of the point and ``size\_color\_glyph'' is the color/glyph/size
information encoded according to internal XGobi specifications
and, therefore, first has to be decoded in ArcView.
\item[XG.RPC.XGobi\_Is\_Up:] Halts ArcView 
while ArcView is waiting for XGobi to start up. 
It does not do anything other than delaying execution for a while.
\item[XG.Select:] Allows the user to select locations in the
current ArcView map view by drawing a boundary based on polygons.
The locations inside the boundary are brushed using the current
selections in the ``Color'', ``Glyph'', and ``Size'' pulldown menus.
Also, remote procedure 16 (described in Section 5.2.1) in XGobi
is called, with the color/glyph/size information of all
selected locations, encoded according to internal XGobi specifications.
\item[XG.Start:] Initiates the ArcView to XGobi link.
It activates the RPC mechanism in ArcView, starts XGobi as a 
background process, redraws the selected view as graphics, and
initializes the new buttons and functionality in the ArcView GUI. In
this script, all global Avenue  variables that are later used in 
other scripts are defined. This script also reads
environment variables and executes the scripts {\bf
RGBColorNameDictionary} and {\bf CreateErrorDictionary}. Internally,
this script calls remote procedures 11, 01, 02, 12, 13, and 14
(described in Section 5.2.1) in XGobi. Each link
type has its own startup script. 
\end{description}

\noindent
Similar Avenue scripts are used to support 
the four other features of the link. 

Since the introduction of the ``Dialog Designer'' extension by ESRI,
we have designed custom dialogs for each of the links. In the
variogram--cloud link, we have introduced a slider to predict the
number of points in the variogram--cloud so that the user can
interactively select the cutoff distance and control the amount of
data passed to XGobi. 


\subsection{Security and Concurrency}

Security concerns the ability of an unauthorized user to ``break'' in
and access or even destroy the ArcView data base.  The
ArcView RPC server supports only one function (function number
``1'': ``script execution -- executes the given script and returns a
string representation of the last object referenced or produced during
the execution of the script'' (Environmental Systems Research
Institute, Inc., 1996)).  
It takes a string as the input where the string must contain
an Avenue script. This script can be
a single Avenue statement or a very long and complex segment of
Avenue code. This could be a potential security problem in the
case that a request by a non--authenticated user
that contains a valid Avenue script of the form
\verb|system.execute("XXX")|, 
where \verb|XXX| may be any Unix command,
is granted.

RPCs and existing security issues have been described in the technical
literature, e.~g., Corbin (1991).  Three levels of authentication
exist. Null authentication means that any RPC request will
be granted. Unix authentication causes the
transmission of additional fields (such as a time stamp, the name of
the local host, and the client's effective user and group IDs) with every
RPC request.  Data Encryption Standard (DES) authentication promotes
secure exchange of data in a standard fashion since it
encrypts/decrypts data through public and private keys associated with
the effective user ID of the calling process. For Unix and DES
authentication, only valid requests are granted.  

ArcView 2.1 and 3.0 provide Unix authentication as the highest
authentication mechanism. In our link between ArcView (2.1 and 3.0) and XGobi
we make use of Unix authentication and process only those
requests where local host and user ID match between ArcView and XGobi
(which is started from within ArcView --- thus IDs will match for
valid requests).  Unfortunately, there is no verifier for Unix
authentication. Thus, the previously mentioned credentials can still
be faked. Therefore, we suggest
that ESRI provides an RPC mechanism for DES authentication in addition
to the current Unix authentication for future releases of
ArcView.

ArcView 2.0 only supported null authentication. 
There was no protection against the manipulation or deletion of entire
ArcView 2.0 data bases or any of the analyst's files by an external
user connected to the running ArcView 2.0 RPC server, hence a total
lack of data security.  

Concurrency means ``What happens if at least two XGobi processes
provide different update information to ArcView at the same time?'' It
has been solved by a built--in feature of XGobi. Only one XGobi
process, the one invoked from within ArcView, provides the features of
an RPC server and RPC client. Any other XGobi processes, e.~g., cloned ones,
are ``regular'' XGobi processes
that do not initiate or respond to RPCs. XGobi processes (the one
invoked from within ArcView and the cloned ones) communicate to each
other through the production and consumption of {\it XEvents}.
Examples of XEvents are the
pressing or releasing of a mouse button, 
crossing a window boundary with the mouse, or events
created by any of the XWindow client programs 
(see Nye, 1992, p.\ 38, for more details).
When points are brushed or subsetted in XGobi, an appropriate XEvent is
generated. XGobi
sequentially processes the XEvents that have been generated by other
XGobi processes. Information will get back to ArcView when the XGobi
invoked from within ArcView processes an XEvent related to brushing or
subsetting.  Due to this
one--to--one link and the serialization of the XEvents, it was not
necessary to consider additional steps to solve the concurrency issue.


\section{Discussion}

One of the current limitations of the ArcView/XGobi link is the number
of points that can be handled. No severe problems arise for the
multivariate and SCDF links since, in both cases, the link reasonably
supports as many points as can be managed in real time within XGobi
itself. However, severe problems arise for
the variogram--cloud link, spatially lagged scatterplot link, and
multivariate variogram--cloud link if the selected cutoff distance
results in hundreds of thousands of points in XGobi.  While XGobi supports
up to about 500,000 points (Swayne et al., 1998), the RPC mechanism
does not. The communication and encoding/decoding of data into strings
that have to be passed between ArcView and XGobi requires considerable
computational time. With large amounts of data, it is possible that
linked brushing between ArcView and XGobi works only partially or
even fails completely when multiple XGobi processes (the one invoked
from within ArcView and the cloned ones) are involved.
For example, when the encoding/decoding of data into strings
takes a larger amount of time for the XGobi process that has been
invoked from within ArcView, this XGobi process might eventually
lose some of the XEvents related to brushing or subsetting information
created by cloned XGobi processes since the XEvent buffer may eventually
be erased after a while. Problems can also occur when a timeout
on the client side makes the client believe that the server does not
work properly any more although the server may be just slow processing
a particular (brushing) request.
There is no immediate solution for these problems due to the restrictions
imposed by the ArcView 3.0 RPC mechanism. It has been indicated by
ESRI that future versions of ArcView will be able to process C
code. This would allow a much faster RPC mechanism where compressed
data is passed between ArcView and XGobi. Both programs might also be
allowed to use shared memory.

Currently, the ArcView/XGobi link does not allow simultaneous
invocations of XGobi with different features. Unfortunately, we cannot
look at an SCDF plot and a variogram--cloud plot at the same time. To
allow such an option in our software environment, we plan to implement
a new type of linked brushing, called ``hierarchical'' linked
brushing. This linked brushing is required for a future environment
where ArcView communicates with several XGobis, each of them
displaying a different feature of the link, e.~g., one XGobi and a
cloned child for the attribute data, one XGobi for the spatial
cumulative distribution functions, and one XGobi for the
variogram--cloud plot. Extensions to our ArcView Avenue code
would be required.

The emphasis of the link is on exploratory
spatial data analysis. ArcView is mostly used to store the data and
display additional geographic features that relate to the statistical
data of interest.  
Currently, the geographic brushing in ArcView exists only
in what Monmonier (1989) defines as its simplest form: ``use a mouse
to highlight specific areas on the map''. 
However, another extension on the ArcView side is
planned to facilitate more complex types of geographic
brushing, e.~g., brushing statements of the form 
``brush all spatial
locations that are at most 10~km away from the next city boundary and
have no major road within a distance of 1~km''. 
This could be easily done now with a query dialog.

An extension of the ArcView/XGobi link towards the statistical
computing environment XploRe (H\"ardle, Klinke \& Turlach, 1995) has
been described in Symanzik, K{\"o}tter, Schmelzer, Klinke, Cook \&
Swayne (1998b), Symanzik, Klinke, Schmelzer, Cook \& Lewin (1997a),
and Symanzik, Cook, Klinke \& Lewin (1998a).  It should be noted that
only few modifications were required on the ArcView/XGobi side to
support this extension while XploRe had to undergo considerable
changes since it did not support any RPC facilities prior to its
addition to the ArcView/XGobi link.  
Recently, an experimental link among the Virtual Reality GIS ViRGIS
(Pajarola, 1998), XGobi, and XploRe has been implemented
(Symanzik, Pajarola \& Widmayer, 1998c) that makes use of the same
RPC technology initially developed for the ArcView/XGobi link.
In the future, other packages
with IPC facilities such as S (Becker, Chambers \& Wilks, 1988;
Chambers, 1998) with its connection class object 
(Chambers, 1997; Chambers, 1998) might be linked as well.
We expect only little difficulties and modifications of the currently
linked three programs when additional programs will be linked.  But
any additional software package most likely requires several changes
to conform to the standard that has been set. However, we do not
expect any major technical problem on either side when including
additional software into the ArcView/XGobi/XploRe/ViRGIS environment.  


\section*{Acknowledgements}

Research related to this article was supported by an EPA grant under
cooperative agreement \#~CR822919--01--0.  The article has not been
subjected to the review of the EPA and thus does not necessarily
reflect the view of the agency and no official endorsement should be
inferred.  Symanzik's research was also partially supported by a
German ``DAAD--Doktorandenstipendium aus Mitteln des zweiten
Hochschulsonderprogramms'' and by a National Science Foundation Group
Infrastructure Grant DMS--9631351. Nicholas Lewin--Koh was supported from a
grant by the John Deere Foundation.  We want to thank Noel Cressie for
the helpful discussion of the theory related to the software. Thanks
are also due to the anonymous reviewers for their comments that helped
in providing an improved presentation of our work.

\section*{Appendix}

The main web address for the software is:
\begin{verbatim}
    http://www.public.iastate.edu/~arcview-xgobi/homepage.html
\end{verbatim}

\noindent
Additional web pages with information are:

\begin{verbatim}
    http://www.math.usu.edu/~symanzik/papers/2000_jcgs/functdesc.ps
    http://www.math.usu.edu/~symanzik/papers/2000_jcgs/AV2XGErrTable
\end{verbatim}


\bibliography{references}


\end{document}

